\chapter{Implementation and reproducibility}
\label{app:impl}

This appendix collects the implementation and reproducibility detail referred to from \chref{methods}.
Code and result artefacts are at \url{https://github.com/W00DSRULES/bio-inspired-foveation-in-dg3}.
Modules named below belong to the \texttt{tez\_deepgaze} package; scripts are named by file.

\section{Dataset access and cross-validation}
\label{app:data}

The dataset is the \pth{MIT1003_initial_fix_consistent} variant of MIT1003, built once by
\texttt{fetch\_mit1003.py --with-initial} and read through \texttt{pysaliency}'s
\pth{get_mit1003_with_initial_fixation}. It returns each scanpath as a variable-length sequence with the
forced central fixation at index~$0$. With \texttt{start\_fixation = 1} that fixation is history only
and every free fixation is scored, the protocol of \citet{kummerer2022deepgaze3} and of \subsecref{mit1003}.
Training and evaluation both select the variant with \texttt{--dataset-variant initial}. The plain
\texttt{get\_mit1003} loader drops the central fixation on load, so the same \texttt{start\_fixation}
then skips the first free fixation as well; it is not used for any number reported here.

Folds are built via \texttt{pysaliency}'s \pth{filter_datasets} module with
\pth{crossval_folds=10}, \pth{val_folds=1} and \pth{test_folds=1}.
The splitter hardcodes \texttt{RandomState(42)} internally, so the shuffle is fixed whatever seed a
caller passes (the \dg{} training notebook passes none), which is why the two
partitions coincide (\subsecref{cv-split}). For each fold $k \in \{0, \ldots, 9\}$ the split has eight
training, one validation, and one test fold.

\section{Evaluator and metrics}
\label{app:eval}

\texttt{evaluate.py:run} accepts any module matching the \dg{} forward signature and never inspects its
class, so one code path produces the pretrained baseline, the sweep tables and the stratified analysis
alike. The per-epoch validation inside the training loop is a separate, lighter loop, pooled per fixation
where the evaluator averages per image; it supplies only the curves of \figref{training-curves}.

NSS is not taken from the \texttt{deepgaze\_pytorch} library. Its released implementation
unpacks \texttt{torch.std\_mean} in the wrong order, so it divides the centred density by the mean where
\eqnref{nss} divides by the standard deviation. \texttt{evaluate.py} therefore carries a local copy with
that line corrected, and every NSS in this thesis comes from it. The pinned dependency set records the
library version but not this substitution, so it is stated here.

The centerbias is the precomputed \texttt{centerbias\_mit1003.npy} from the \dg{} release. At evaluation
time \texttt{centerbias.py} resizes it to each stimulus by nearest-neighbour zoom and renormalises in
log-space via \texttt{logsumexp}, so the result is a proper log-probability density summing to one over
image pixels.

\section{Consensus regions}
\label{app:consensus}

The figures of \chref{results} draw inter-subject agreement as a \emph{consensus region}: the part of
the image that at least three-quarters of the subjects fixated within $2^\circ$ of, enclosed by a
contour, with a second contour for the $90\%$ core. Two degrees is the foveal radius of
\secref{foveated-vision}, so a pixel inside the region is one those subjects saw foveally. The region is
a way of drawing agreement rather than a measure of it; the scalar the correlations of \chref{results}
read is the entropy of \eqnref{fix-entropy}.

Agreement is summarised by the \emph{consensus count map}, an array
the shape of the image whose value at each pixel is the number of distinct subjects who fixated within
a radius $R$ of it. For a stimulus of shape $(H, W)$, $N$ subjects, and radius $R$,

\begin{equation}
  C(y, x) \;=\; \sum_{s=1}^{N} \mathbf{1}\!\left[
    \exists\, (x_i, y_i) \in \mathcal{F}_s \;:\; \norm{(y, x) - (y_i, x_i)} \leq R
  \right],
  \label{eq:consensus-count}
\end{equation}

where $\mathcal{F}_s$ is subject $s$'s fixation set and $\mathbf{1}[\cdot]$ the indicator function.
Thresholding at fixed fractions of the subject count gives the two contours the figures draw:

\begin{equation}
  M_{0.75}(y, x) = \mathbf{1}[C(y, x) \geq \lceil 0.75 N \rceil],
  \qquad
  M_{0.90}(y, x) = \mathbf{1}[C(y, x) \geq \lceil 0.90 N \rceil].
  \label{eq:consensus-masks}
\end{equation}

$N$ counts the subjects who fixated the image, so with fifteen subjects the two thresholds resolve to
at least twelve and at least fourteen of them. The radius is $R = 70$~pixels throughout, which is
$2^\circ$ at the MIT1003 viewing geometry.

\section{Input-foveation implementation}
\label{app:foveation}

\subsection*{Where the human cutoff \texorpdfstring{$f_{c0} = 40$}{fc0 = 40} comes from}

\subsecref{gp-model} takes $f_{c0} = 40$~cyc/deg as human acuity rather than the ${\sim}50$~cyc/deg
\citet{geislerperry1998} quote in their introduction. The two figures come from different places in the
same paper: the $50$ is a general statement of human resolution, paired with a half-acuity eccentricity of
$2.5^\circ$, whereas the model they fitted is a contrast-threshold law,

\begin{equation}
  CT(f, e) \;=\; CT_0 \exp\!\left(\alpha f \, \frac{e + e_2}{e_2}\right).
  \label{eq:gp-contrast}
\end{equation}

$CT$ is the faintest contrast at which a grating of spatial frequency $f$ is still visible at
eccentricity $e$. Of its three fitted constants, $CT_0$ is that threshold at the fovea for a coarse
grating, $\alpha$ sets how fast the threshold rises as the grating gets finer, and $e_2$ how fast it
rises away from the fovea.

The bracket $(e + e_2)/e_2$ is the eccentricity penalty, equal to one at the fovea and doubling at
$e = e_2$. At the fovea it drops out, leaving $CT(f, 0) = CT_0 \exp(\alpha f)$.

A cutoff is a frequency above which nothing is visible however the image is drawn, which is where the
threshold reaches $CT = 1$, the maximum contrast. Setting $CT = 1$ in the foveal case and solving for
$f$ leaves the foveal cutoff fixed by two fitted constants alone:

\begin{equation}
  f_{c0} \;=\; \frac{\ln(1/CT_0)}{\alpha}.
  \label{eq:fc0-derivation}
\end{equation}

Solving the same way at general $e$, keeping the bracket, returns \eqnref{gp-acuity}, so the acuity law
\secref{foveation} is built on is this fitted law rearranged rather than a second assumption.

\citet{geislerperry1998} fit \eqnref{gp-contrast} to the contrast-sensitivity measurements of
\citet{robson1981}, obtaining $\alpha = 0.106$, $e_2 = 2.3^\circ$ and $CT_0 = 1/64$, and report that the
same $\alpha$ and $e_2$ also fit two further published contrast-sensitivity studies with a different
minimum contrast threshold each: $CT_0 = 1/75$ for \citet{arnow1996} and $1/76$ for \citet{banks1991}.
\eqnref{fc0-derivation} turns those three thresholds into $39.2$, $40.7$ and $40.9$~cyc/deg
respectively, a mean of $40.3$ with a standard deviation of $0.9$. That spread is a check on how much the foveal cutoff moves when the threshold is
measured on different subjects, not three independent refits: $\alpha$ and $e_2$ are held fixed across
all three, and only $CT_0$ varies. The $f_{c0} = 40$ used here and the $e_2 = 2.3^\circ$ of
\subsecref{gp-model} therefore come from the same fit.

\subsection*{Viewing geometry and the pixels-per-degree constant}

\subsecref{gp-model} uses $\ppd = 35$ for MIT1003. \citet{judd2009mit1003} report viewing
\enquote{approximately two feet from a 19 inch computer screen of resolution $1280 \times 1024$}, which
works out to $\ppd \approx 36.1$. The \dg{} release states $35$~pixels per degree for the same dataset
\citep{deepgazepytorch}, which is the value used here.
In pixel space $\ppd$ enters the blur only as an additive offset: substituting
$e_{\text{deg}} = e_{\text{px}}/\ppd$ into \eqnref{gp-acuity} and \eqnref{frac-level} gives

\begin{equation}
  L(e_{\text{px}}) \;=\; \log_2 \frac{e_{\text{px}} + e_2\,\ppd}{2 f_{c0} e_2},
  \label{eq:level-pixels}
\end{equation}

so the far field does not depend on $\ppd$ at all. By \eqnref{level-pixels} the difference between $35$
and $36.1$ shifts the blur by at most $3\%$ in $\sigma$, and by nothing inside the sharp disc, where the
clamp holds $L = 0$.

\subsection*{Implementation}

The space-variant foveation of \secref{foveation} is implemented in \texttt{foveate\_input.py}. The
Geisler--Perry cutoff \eqnref{gp-acuity} is evaluated per pixel from the eccentricity to the current
fixation and the pixels-per-degree constant, converted to cycles-per-pixel, and clamped at the display
Nyquist ($0.5$~cyc/px); the fractional pyramid level \eqnref{frac-level} then indexes a Gaussian pyramid
whose two bracketing levels are blended with triangular weights that sum to one. The pyramid is octave
spaced: level $i$ targets a half-amplitude cutoff of $0.5 \cdot 2^{-i}$~cyc/px, which for a Gaussian of
standard deviation $\sigma$ gives $\sigma_i = 0.375 \cdot 2^{i}$~px, and the levels are built with
\texttt{torchvision}'s differentiable \texttt{gaussian\_blur}. The transform is geometry-preserving
(pixels do not move), so the output is never a log-polar warp. In the gaze-contingent arms the image is
foveated around the most recent fixation of the history before each forward pass; the pyramid is built
once per image and reused across the fixations scored on it.

\section{Training internals}
\label{app:training}

Read-out fine-tuning is implemented in \texttt{foveated\_train.py}. \pth{freeze_for_head_training} sets
\texttt{requires\_grad=False} on the DenseNet backbone (\texttt{features}) and the finalizers and places
every submodule in \texttt{.eval()} mode, so BatchNorm running statistics and the finalizer's two learned
parameters, its Gaussian kernel width and its centerbias weight, stay fixed; the saliency, scanpath, and
fixation-selection networks remain trainable.
\texttt{audit\_trainable} asserts that every trainable parameter lies in those three read-out groups and that
the frozen groups contribute none, raising otherwise. The arms differ only in the input handed to the
same training routine: sharp, foveated, or foveated with the centre pinned to the image midpoint for
every fixation.

The optimiser is Adam over the read-out parameters at a base learning rate
of $3 \times 10^{-4}$; \texttt{lr\_for\_epoch} applies the step decay of \subsecref{training-loop},
dividing the base rate by ten at epoch five. Its milestones are absolute, at epochs five, eight and
eleven, so the two later ones fall outside the seven-epoch schedule and never fire.

One image yields about a hundred (history, target) pairs, so the per-image loss is accumulated over
\texttt{micro\_batch} chunks before one optimiser step, and gradients are clipped to unit norm.

Every epoch writes one checkpoint bundle: \texttt{weights.pt}, which holds only the trainable read-out
tensors, since the frozen submodules stay byte-identical to the pretrained weights at every epoch,
\texttt{optimizer.pt} with Adam's moment estimates, and a \texttt{metrics.json} sidecar. The sidecar
records the validation log-likelihood and information gain, the seed, the git commit, and
SHA-256 hashes of \texttt{uv.lock} (which pins every dependency, the \texttt{deepgaze-pytorch} git
revision and the \texttt{pysaliency} release among them), the \texttt{deepgaze3.pth} weights, the
\texttt{centerbias\_mit1003.npy} prior and the cross-validation split file, so a changed input shows as a
hash mismatch.

\section{Cluster execution}
\label{app:cluster}

\texttt{device.py:pick\_device} chooses CUDA, then MPS, then CPU. Every reported run used CUDA on the
NVIDIA \texttt{gpu2} partition of the Goethe-NHR cluster, one Slurm job per arm and fold, from seed $42$.
The driver \pth{foveation_pump.sh} submits the seven-arm matrix of \subsecref{arms}
(sharp, gaze-contingent at $f_{c0} \in \{40, 20, 10\}$ and fixed-centre at the same three
cutoffs, ten folds each, so $70$ jobs). Training was launched as
\begin{center}
\texttt{TEZ\_EPOCHS=7 TEZ\_DATASET\_VARIANT=initial bash} \pth{scripts/slurm/foveation_pump.sh}
\end{center}
and the tables of \chref{results} come from \pth{foveation_sweep_eval.sbatch}, run once per split with
\texttt{TEZ\_EPOCH=5} and \texttt{TEZ\_DATASET\_VARIANT=initial} and scoring the checkpoints with the
evaluator of \secref{eval-framework}; \pth{foveation_sweep_table.py} assembles
\texttt{table.json} and \pth{make_results_figures.py} renders \tabref{phase17-arms} from it. The stratified
diagnostics come from \pth{foveation_stratified.py} at the same epoch on the validation split, and the
per-image diagnostic from \texttt{per\_image\_diagnostic.py --n-stim 1003 --dataset-variant initial}.
